\documentclass[11pt]{article}
\usepackage[paperwidth=220mm,paperheight=320mm,left=16mm,right=16mm,top=18mm,bottom=18mm,headheight=16pt,headsep=10pt,footskip=14pt]{geometry}
\usepackage{fontspec,graphicx,tikz}
\usepackage[unicode,hidelinks]{hyperref}
\usepackage{polyglossia}
\setmainlanguage{latin}
\setotherlanguages{arabic}
\setmainfont{Linux Libertine O}[Ligatures=TeX]
\newfontfamily\arabicfont{Amiri}[Script=Arabic]
\newfontfamily\ArBody{Amiri}[Script=Arabic]
\newfontfamily\OrnFont{FreeSerif}
\hypersetup{pdftitle={Al-Battani, Kitab az-Zij as-Sabi'. Textus Arabicus (Nallino, Pars III). S10},pdfauthor={Al-Battani; Carlo Alfonso Nallino (ed.)}}
\makeatletter
\def\ps@sourceedition{%
\def\@oddhead{\vbox{\hbox to\textwidth{\fontsize{8}{10}\selectfont AL-BATTĀNĪ — TEXTUS ARABICUS (NALLINO, PARS III)\hfil\leftmark}\vskip3pt\hrule height.25pt}}%
\let\@evenhead\@oddhead
\def\@oddfoot{\hbox to\textwidth{\fontsize{7}{9}\selectfont\rightmark\hfil\thepage}}%
\let\@evenfoot\@oddfoot}
\makeatother
\pagestyle{sourceedition}
\setlength{\parindent}{0pt}\setlength{\parskip}{0pt}
\tracinglostchars=2
\newcommand{\DSource}[2]{\clearpage\markboth{#2}{AB01-PDF#1}\pdfbookmark[0]{AB01-PDF#1 — #2}{bk-#1}\null}
% place #3 with its anchor #1 at (#2) from the top left corner of the page
\newcommand{\At}[3]{\begin{tikzpicture}[remember picture,overlay]\node[anchor=#1,inner sep=0pt,outer sep=0pt] at ([shift={(#2)}]current page.north west) {#3};\end{tikzpicture}}
\newsavebox{\LBox}
% a right-to-left line #2 justified to the width #1 (\RLJ) or at its natural width (\RLN); narrowed to #1 if wider
\newcommand{\RLJ}[2]{\sbox{\LBox}{\RL{#2}}%
 \ifdim\wd\LBox>#1\relax\typeout{NARROWED|\the\wd\LBox|#1}\resizebox{#1}{\ht\LBox}{\usebox{\LBox}}%
 \else\makebox[#1][s]{\RL{#2}}\fi}
\newcommand{\RLN}[2]{\sbox{\LBox}{\RL{#2}}%
 \ifdim\wd\LBox>#1\relax\typeout{NARROWED|\the\wd\LBox|#1}\resizebox{#1}{\ht\LBox}{\usebox{\LBox}}%
 \else\usebox{\LBox}\fi}
\newcommand{\Note}[1]{\textsuperscript{\fontsize{8}{8}\selectfont\addfontfeatures{Numbers=Lining}#1}}
\newcommand{\Ov}[1]{\vbox{\hrule height .4pt\kern 1.0pt\hbox{#1}}}
\newcommand{\ArRun}[1]{\textarabic{#1}}
\begin{document}
\thispagestyle{empty}
\begin{center}\Large AL-BATTĀNĪ\par\vspace{4mm}\LARGE KITĀB AZ-ZĪJ AṢ-ṢĀBIʾ\par
\vspace{5mm}\large textus Arabicus a Carolo Alphonso Nallino editus (pars III)\par\vspace{8mm}
\Large S10 --- editio diplomatica\par\vspace{4mm}\large Titulus; paginae impressae 1--7\end{center}
\vspace{12mm}\noindent Singulae lineae paginarum impressarum singulis lineis huius editionis respondent et ibi
collocantur ubi in pagina impressa stant (lineae quae totam latitudinem implent ad eandem latitudinem
extenduntur). Textus Arabicus cum signis vocalium, ut impressa sunt, ex imaginibus paginarum transcriptus est;
productiones litterarum (kashida) non transcribuntur. Numeri foliorum codicis et numeri linearum Nallini in
marginibus stant ubi impressi sunt. Haec transcriptio a nullo homine recognita est.\par
\DSource{1153}{SINE NUMERO}
\At{base}{326.68pt,-145.19pt}{\hypertarget{AB01-PDF1153-L01}{}{\fontsize{46.0}{55.2}\ArBody\RL{كتاب الزيج الصابئ}}}
\At{base}{324.13pt,-215.19pt}{\hypertarget{AB01-PDF1153-L02}{}{\fontsize{13.0}{15.6}\ArBody\RL{تاليف}}}
\At{base}{325.58pt,-258.88pt}{\hypertarget{AB01-PDF1153-L03}{}{\fontsize{27.0}{32.4}\ArBody\RL{ابي عبد الله محمّد بن سنان بن جابر الحرّانيّ}}}
\At{base}{324.73pt,-318.69pt}{\hypertarget{AB01-PDF1153-L04}{}{\fontsize{27.0}{32.4}\ArBody\RL{المعروف بالبَتّانيّ}}}
\At{base}{325.23pt,-384.58pt}{\hypertarget{AB01-PDF1153-L05}{}{\fontsize{19.0}{22.8}\ArBody\RL{نُقِلَ عن النسخة المحفوظة بمكتبة بلدة الإسْكوريال من بلاد الاندلس}}}
\At{base}{324.18pt,-442.19pt}{\hypertarget{AB01-PDF1153-L06}{}{\fontsize{11.5}{13.8}\ArBody\RL{اعتنى بطبعه وتصحيحه وترجمه الى اللغة اللاتينيّة وعلّق حواشيه}}}
\At{base}{323.08pt,-482.49pt}{\hypertarget{AB01-PDF1153-L07}{}{\fontsize{17.0}{20.4}\ArBody\RL{الدكتور كَرْلو نالّينو}}}
\At{base}{324.53pt,-516.99pt}{\hypertarget{AB01-PDF1153-L08}{}{\fontsize{11.5}{13.8}\ArBody\RL{مدرّس بمدرسة اللغات الشرقيّة في نابولي}}}
\At{center}{325.58pt,-604.48pt}{{\OrnFont\fontsize{20}{20}\selectfont ❦}}
\At{base}{323.63pt,-708.88pt}{\hypertarget{AB01-PDF1153-L09}{}{\fontsize{22.0}{26.4}\ArBody\RL{طُبع بمدينة رومية العظمى}}}
\At{base}{322.93pt,-741.29pt}{\hypertarget{AB01-PDF1153-L10}{}{\fontsize{11.5}{13.8}\ArBody\RL{سنة ١٨٩٩ المسيحيّة}}}
\DSource{1151}{PAG. 1}
\At{base}{338.23pt,-147.44pt}{\hypertarget{AB01-PDF1151-L01}{}{\fontsize{26.0}{31.2}\ArBody\RL{بسم الله الرحمن الرحيم}}}
\At{base east}{115.03pt,-147.44pt}{{\fontsize{8}{9}\selectfont fol. 1,v.}}
\At{base}{333.18pt,-214.75pt}{\hypertarget{AB01-PDF1151-L02}{}{\fontsize{26.0}{31.2}\ArBody\RL{صلّى الله على النبيّ محمّد رسوله الكريم وعلى آله وصحبه وسلّم}}}
\At{center}{331.33pt,-276.84pt}{\rule{69.0pt}{0.5pt}}
\At{base east}{521.43pt,-411.64pt}{\hypertarget{AB01-PDF1151-L03}{}{\fontsize{13.8}{27}\ArBody\RLJ{401.40pt}{جامع ما وضع محمّد بن جابر بن سنان الحرّانيّ المعروف بالبَتّانيّ عفا الله عنه في حساب النجوم}}}
\At{base east}{545.53pt,-438.64pt}{\hypertarget{AB01-PDF1151-L04}{}{\fontsize{13.8}{27}\ArBody\RLN{358.50pt}{ومواضع مسيرها الممتحن وجملة ما فيه من الأبواب سبعة\Note{1} وخمسون بابًا وهذا تفسيرها.}}}
\At{base east}{520.73pt,-495.14pt}{\hypertarget{AB01-PDF1151-L05}{}{\fontsize{13.8}{27}\ArBody\RLN{232.90pt}{في صدر الكتاب}}}
\At{base east}{545.53pt,-495.14pt}{{\fontsize{13.8}{27}\ArBody\Ov{ا}}}
\At{base east}{520.03pt,-518.94pt}{\hypertarget{AB01-PDF1151-L06}{}{\fontsize{13.8}{27}\ArBody\RLN{385.20pt}{في تقسيم دائرة الفلك وضرب الأجزاء بعضها في بعض وتجذيرها وقسمتها بعضها على بعض}}}
\At{base east}{560.33pt,-518.94pt}{{\fontsize{13.8}{27}\ArBody\Ov{ب}}}
\At{base west}{555.53pt,-518.94pt}{{\fontsize{8}{9}\selectfont 15}}
\At{base east}{519.63pt,-547.35pt}{\hypertarget{AB01-PDF1151-L07}{}{\fontsize{13.8}{27}\ArBody\RLJ{400.70pt}{في معرفة اقدار اوتار اجزاء الدائرة وإِثبات أنصاف اوتار أضعاف القسيّ في الجداول وما يتبع}}}
\At{base east}{545.53pt,-547.35pt}{{\fontsize{13.8}{27}\ArBody\Ov{ج}}}
\At{base east}{520.73pt,-576.14pt}{\hypertarget{AB01-PDF1151-L08}{}{\fontsize{13.8}{27}\ArBody\RLN{81.00pt}{ذلك من العمل بها}}}
\At{base east}{520.03pt,-599.55pt}{\hypertarget{AB01-PDF1151-L09}{}{\fontsize{13.8}{27}\ArBody\RLJ{401.40pt}{في مقدار ميل فلك البروج عن فلك معدّل النهار ونجزية هذا الميل وجهاته ومراتبه في صعوده}}}
\At{base east}{545.53pt,-599.55pt}{{\fontsize{13.8}{27}\ArBody\Ov{د}}}
\At{base east}{520.03pt,-629.14pt}{\hypertarget{AB01-PDF1151-L10}{}{\fontsize{13.8}{27}\ArBody\RLN{190.10pt}{وهبوطه وهو ميل الشمس عن الفلك المستقيم}}}
\At{base east}{519.33pt,-654.64pt}{\hypertarget{AB01-PDF1151-L11}{}{\fontsize{13.8}{27}\ArBody\RLJ{402.20pt}{في معرفة اقدار ما يطلع من فلك معدّل النهار مع قسيّ فلك البروج المفروضة\Note{2} تحت معدّل النهار}}}
\At{base east}{559.23pt,-654.64pt}{{\fontsize{13.8}{27}\ArBody\Ov{ه}}}
\At{base west}{555.53pt,-654.64pt}{{\fontsize{8}{9}\selectfont 20}}
\At{base east}{519.33pt,-682.05pt}{\hypertarget{AB01-PDF1151-L12}{}{\fontsize{13.8}{27}\ArBody\RLJ{401.80pt}{الذي يسمَّى خطّ الاستواء وبهذه الاقدار ايضًا تمرّ البروج وتجوز في فلك نصف النهار في كلّ}}}
\At{base east}{518.93pt,-708.64pt}{\hypertarget{AB01-PDF1151-L13}{}{\fontsize{13.8}{27}\ArBody\RLN{272.20pt}{موضع من الارض ويسمَّى لذلك مطالع البروج في الفلك المستقيم}}}
\At{west}{117.53pt,-735.14pt}{\rule{102.2pt}{0.5pt}}
\At{base west}{143.43pt,-763.55pt}{\hypertarget{AB01-PDF1151-N01}{}{\fontsize{9.2}{11}\selectfont 1) Cod. \ArRun{سبع} — 2) Excidisse videtur \ArRun{في البلدان التي}}}
\DSource{1150}{PAG. 2}
\At{base east}{498.63pt,-143.27pt}{\hypertarget{AB01-PDF1150-L01}{}{\fontsize{13.8}{27}\ArBody\RLJ{397.80pt}{في معرفة خواصّ كلّ خطّ من الخطوط الموازية لمعدّل النهار المائل عنه الى الشمال وذكر مواضع}}}
\At{base east}{524.73pt,-143.27pt}{{\fontsize{13.8}{27}\ArBody\Ov{و}}}
\At{base east}{497.93pt,-170.27pt}{\hypertarget{AB01-PDF1150-L02}{}{\fontsize{13.8}{27}\ArBody\RLN{265.70pt}{الارض العامرة المعلومة الطول والعرض في كتاب صورة الارض}}}
\At{base east}{498.63pt,-197.27pt}{\hypertarget{AB01-PDF1150-L03}{}{\fontsize{13.8}{27}\ArBody\RLJ{397.80pt}{في معرفة سعة مشارق الشتاء والصيف ومغاربها من دوائر آفاق البلدان وهي القسيّ التي تكون}}}
\At{base east}{524.33pt,-197.27pt}{{\fontsize{13.8}{27}\ArBody\Ov{ز}}}
\At{base east}{498.63pt,-224.27pt}{\hypertarget{AB01-PDF1150-L04}{}{\fontsize{13.8}{27}\ArBody\RLJ{397.80pt}{بين فلك معدّل النهار ومواضع فلك البروج في دائرة الافق ويسمَّى سمت المطالع والمغارب من}}}
\At{base east}{498.63pt,-251.27pt}{\hypertarget{AB01-PDF1150-L05}{}{\fontsize{13.8}{27}\ArBody\RLN{45.00pt}{دائرة الافق}}}
\At{base east}{91.63pt,-251.27pt}{{\fontsize{8}{9}\selectfont 5}}
\At{base east}{498.63pt,-278.56pt}{\hypertarget{AB01-PDF1150-L06}{}{\fontsize{13.8}{27}\ArBody\RLN{324.00pt}{في معرفة ارتفاع القطب الشماليّ من قبل زيادة النهار الاطول اذا كان مفروضًا}}}
\At{base east}{525.03pt,-278.56pt}{{\fontsize{13.8}{27}\ArBody\Ov{ح}}}
\At{base east}{498.63pt,-305.97pt}{\hypertarget{AB01-PDF1150-L07}{}{\fontsize{13.8}{27}\ArBody\RLN{267.50pt}{في معرفة زيادة *النهار الاطول من قبل ارتفاع القطب المفروض}}}
\At{base east}{525.03pt,-305.97pt}{{\fontsize{13.8}{27}\ArBody\Ov{ط}}}
\At{base west}{531.53pt,-305.97pt}{{\fontsize{8}{9}\selectfont f. 2,r.}}
\At{base east}{498.63pt,-332.56pt}{\hypertarget{AB01-PDF1150-L08}{}{\fontsize{13.8}{27}\ArBody\RLJ{397.80pt}{في معرفة الارتفاع والظلّ احدهما من قبل الآخر اذا كان الظلّ بسيطًا ومعرفة ذلك اذا كان}}}
\At{base east}{525.43pt,-332.56pt}{{\fontsize{13.8}{27}\ArBody\Ov{ى}}}
\At{base east}{497.93pt,-359.56pt}{\hypertarget{AB01-PDF1150-L09}{}{\fontsize{13.8}{27}\ArBody\RLN{40.40pt}{الظلّ قائمًا}}}
\At{base east}{498.63pt,-386.56pt}{\hypertarget{AB01-PDF1150-L10}{}{\fontsize{13.8}{27}\ArBody\RLJ{397.80pt}{في معرفة سمت الارتفاع والظلّ من دائرة الافق في كلّ بلد وفي كلّ وقت من النهار في جميع}}}
\At{base east}{525.03pt,-386.56pt}{{\fontsize{13.8}{27}\ArBody\Ov{يا}}}
\At{base east}{91.93pt,-386.56pt}{{\fontsize{8}{9}\selectfont 10}}
\At{base east}{498.23pt,-413.56pt}{\hypertarget{AB01-PDF1150-L11}{}{\fontsize{13.8}{27}\ArBody\RLJ{397.40pt}{اجزاء فلك البروج وهو ما تقطع القوس التي تجوز على سمت الرؤس والشمس من دائرة الافق}}}
\At{base east}{498.63pt,-440.97pt}{\hypertarget{AB01-PDF1150-L12}{}{\fontsize{13.8}{27}\ArBody\RLN{90.40pt}{من حدّ المطلع والمغيب}}}
\At{base east}{498.63pt,-467.97pt}{\hypertarget{AB01-PDF1150-L13}{}{\fontsize{13.8}{27}\ArBody\RLJ{397.80pt}{في معرفة خطّ نصف النهار في كلّ بلد وهو سمت الجنوب وما يظهر معه من سمت مشرق}}}
\At{base east}{524.73pt,-467.97pt}{{\fontsize{13.8}{27}\ArBody\Ov{يب}}}
\At{base east}{498.63pt,-494.56pt}{\hypertarget{AB01-PDF1150-L14}{}{\fontsize{13.8}{27}\ArBody\RLN{122.40pt}{الاعتدال ومعرفته بجهات شتَّى}}}
\At{base east}{498.63pt,-521.26pt}{\hypertarget{AB01-PDF1150-L15}{}{\fontsize{13.8}{27}\ArBody\RLJ{397.80pt}{في معرفة قدر ما يطلع من فلك معدّل النهار مع اجزاء فلك البروج المفروضة من الافق في كلّ}}}
\At{base east}{525.43pt,-521.26pt}{{\fontsize{13.8}{27}\ArBody\Ov{يج}}}
\At{base east}{92.33pt,-521.26pt}{{\fontsize{8}{9}\selectfont 15}}
\At{base east}{498.63pt,-548.57pt}{\hypertarget{AB01-PDF1150-L16}{}{\fontsize{13.8}{27}\ArBody\RLJ{397.80pt}{موضع من مواضع الارض ويسمَّى مطالع البروج في كلّ بلد وما يتبع ذلك من معرفة مطالع ايّ}}}
\At{base east}{498.63pt,-575.57pt}{\hypertarget{AB01-PDF1150-L17}{}{\fontsize{13.8}{27}\ArBody\RLJ{397.80pt}{وجه شئت في هذه المطالع وفي مطالع الفلك المستقيم ومعرفة اجزاء فلك البروج من قبل هذه}}}
\At{base east}{498.63pt,-602.57pt}{\hypertarget{AB01-PDF1150-L18}{}{\fontsize{13.8}{27}\ArBody\RLJ{397.80pt}{المطالع ومقدار قوس النهار والليل وساعاتها المعتدلة وازمان ساعات النهار والليل الزمانيّة وتحويل}}}
\At{base east}{498.63pt,-629.57pt}{\hypertarget{AB01-PDF1150-L19}{}{\fontsize{13.8}{27}\ArBody\RLN{63.70pt}{بعضها الى بعض}}}
\At{base east}{498.63pt,-656.26pt}{\hypertarget{AB01-PDF1150-L20}{}{\fontsize{13.8}{27}\ArBody\RLN{307.80pt}{في معرفة عروض البلدان وهو ارتفاع القطب الشماليّ بها عن الافق بالرصد}}}
\At{base east}{524.73pt,-656.26pt}{{\fontsize{13.8}{27}\ArBody\Ov{يد}}}
\At{base east}{92.33pt,-656.26pt}{{\fontsize{8}{9}\selectfont 20}}
\At{base east}{498.63pt,-682.87pt}{\hypertarget{AB01-PDF1150-L21}{}{\fontsize{13.8}{27}\ArBody\RLN{245.50pt}{في معرفة ارتفاع الشمس في وقت إنصاف النهار في كلّ يوم}}}
\At{base east}{524.73pt,-682.87pt}{{\fontsize{13.8}{27}\ArBody\Ov{يه}}}
\At{base east}{498.63pt,-709.87pt}{\hypertarget{AB01-PDF1150-L22}{}{\fontsize{13.8}{27}\ArBody\RLJ{397.80pt}{في معرفة ما يمضي من النهار من ساعة وما يطلع من قبل قياس الشمس ومعرفة الارتفاع والظلّ}}}
\At{base east}{524.73pt,-709.87pt}{{\fontsize{13.8}{27}\ArBody\Ov{يو}}}
\At{base east}{497.53pt,-736.87pt}{\hypertarget{AB01-PDF1150-L23}{}{\fontsize{13.8}{27}\ArBody\RLN{139.70pt}{القائم}}}
\At{base east}{498.63pt,-763.87pt}{\hypertarget{AB01-PDF1150-L24}{}{\fontsize{13.8}{27}\ArBody\RLN{227.20pt}{في معرفة الارتفاع من قبل ما يمضي من ساعات النهار}}}
\At{base east}{524.33pt,-763.87pt}{{\fontsize{13.8}{27}\ArBody\Ov{يز}}}
\DSource{1149}{PAG. 3}
\At{base east}{522.03pt,-143.44pt}{\hypertarget{AB01-PDF1149-L01}{}{\fontsize{13.8}{27}\ArBody\RLJ{399.60pt}{في معرفة ابعاد الكواكب الثابتة او المتحيّرة عن فلك معدّل النهار اذا كانت مائلةً عن نطاق البروج}}}
\At{base east}{544.83pt,-143.44pt}{{\fontsize{13.8}{27}\ArBody\Ov{يج}}}
\At{base east}{522.03pt,-170.44pt}{\hypertarget{AB01-PDF1149-L02}{}{\fontsize{13.8}{27}\ArBody\RLJ{399.60pt}{في العرض واجزاء فلك البروج التي تتوسّط السماء معها من قبل مواضعها من *فلك البروج في}}}
\At{base east}{113.23pt,-170.44pt}{{\fontsize{8}{9}\selectfont f. 2,v.}}
\At{base east}{522.03pt,-197.44pt}{\hypertarget{AB01-PDF1149-L03}{}{\fontsize{13.8}{27}\ArBody\RLN{63.00pt}{الطول والعرض}}}
\At{base east}{522.03pt,-224.44pt}{\hypertarget{AB01-PDF1149-L04}{}{\fontsize{13.8}{27}\ArBody\RLJ{399.60pt}{في معرفة نصف\Note{1} قوس نهار احد الكواكب وهو نصف مكثه فوق الارض وتحتها ايضًا وازمان}}}
\At{base east}{545.93pt,-224.44pt}{{\fontsize{13.8}{27}\ArBody\Ov{يط}}}
\At{base east}{522.03pt,-250.34pt}{\hypertarget{AB01-PDF1149-L05}{}{\fontsize{13.8}{27}\ArBody\RLN{110.90pt}{ساعاته فوق الارض وتحتها}}}
\At{base west}{552.43pt,-250.34pt}{{\fontsize{8}{9}\selectfont 5}}
\At{base east}{522.03pt,-278.14pt}{\hypertarget{AB01-PDF1149-L06}{}{\fontsize{13.8}{27}\ArBody\RLN{374.10pt}{في معرفة الدرجة من فلك البروج التي يطلع معها احد الكواكب والدرجة التي معها يغيب}}}
\At{base east}{545.93pt,-278.14pt}{{\fontsize{13.8}{27}\ArBody\Ov{ك}}}
\At{base east}{522.03pt,-305.14pt}{\hypertarget{AB01-PDF1149-L07}{}{\fontsize{13.8}{27}\ArBody\RLN{259.90pt}{في معرفة ما يمضي من الليل من ساعة بقياس بعض الكواكب}}}
\At{base east}{546.63pt,-305.14pt}{{\fontsize{13.8}{27}\ArBody\Ov{كا}}}
\At{base east}{522.03pt,-331.75pt}{\hypertarget{AB01-PDF1149-L08}{}{\fontsize{13.8}{27}\ArBody\RLN{290.20pt}{في معرفة ارتفاع بعض الكواكب من قبل الساعات الماضية من الليل}}}
\At{base east}{548.83pt,-331.75pt}{{\fontsize{13.8}{27}\ArBody\Ov{كب}}}
\At{base east}{522.03pt,-358.75pt}{\hypertarget{AB01-PDF1149-L09}{}{\fontsize{13.8}{27}\ArBody\RLN{244.80pt}{في معرفة سمت احد الكواكب من قبل ارتفاعه عن الافق}}}
\At{base east}{547.03pt,-358.75pt}{{\fontsize{13.8}{27}\ArBody\Ov{كج}}}
\At{base east}{522.03pt,-385.75pt}{\hypertarget{AB01-PDF1149-L10}{}{\fontsize{13.8}{27}\ArBody\RLJ{399.60pt}{في معرفة بعد احد الكواكب عن فلك معدّل النهار وما يتوسّط السماء معه من اجزاء البروج}}}
\At{base east}{548.83pt,-385.75pt}{{\fontsize{13.8}{27}\ArBody\Ov{كد}}}
\At{base west}{555.33pt,-385.75pt}{{\fontsize{8}{9}\selectfont 10}}
\At{base east}{522.03pt,-412.75pt}{\hypertarget{AB01-PDF1149-L11}{}{\fontsize{13.8}{27}\ArBody\RLJ{399.60pt}{من قبل معرفة سمت الموضع الذي يطلع منه او يغيب من دائرة الافق. وبه يُعْلَم ايضًا ميل الجزء}}}
\At{base east}{522.03pt,-439.35pt}{\hypertarget{AB01-PDF1149-L12}{}{\fontsize{13.8}{27}\ArBody\RLN{164.20pt}{من فلك البروج عن فلك معدّل النهار}}}
\At{base east}{522.03pt,-466.75pt}{\hypertarget{AB01-PDF1149-L13}{}{\fontsize{13.8}{27}\ArBody\RLJ{399.60pt}{في معرفة الجزء الذي فيه الكوكب من اجزاء فلك البروج وعرض الكوكب من قبل بعده عن}}}
\At{base east}{549.53pt,-466.75pt}{{\fontsize{13.8}{27}\ArBody\Ov{كه}}}
\At{base east}{522.03pt,-493.35pt}{\hypertarget{AB01-PDF1149-L14}{}{\fontsize{13.8}{27}\ArBody\RLN{274.00pt}{فلك معدّل النهار والجزء الذي يتوسّط السماء معه اذا كان معلومًا}}}
\At{base east}{522.03pt,-520.35pt}{\hypertarget{AB01-PDF1149-L15}{}{\fontsize{13.8}{27}\ArBody\RLN{348.90pt}{في معرفة ابعاد ما بين الكواكب على ترتيب مواضعها في الفلك في الطول والعرض}}}
\At{base east}{548.13pt,-520.35pt}{{\fontsize{13.8}{27}\ArBody\Ov{كو}}}
\At{base west}{555.63pt,-520.35pt}{{\fontsize{8}{9}\selectfont 15}}
\At{base east}{522.03pt,-547.04pt}{\hypertarget{AB01-PDF1149-L16}{}{\fontsize{13.8}{27}\ArBody\RLJ{399.60pt}{في معرفة مقدار طول ازمان\Note{2} السنة الشمسيّة الموجودة بالرصد وحركة الشمس الوسطى في الايّام}}}
\At{base east}{549.13pt,-547.04pt}{{\fontsize{13.8}{27}\ArBody\Ov{كز}}}
\At{base east}{522.03pt,-573.64pt}{\hypertarget{AB01-PDF1149-L17}{}{\fontsize{13.8}{27}\ArBody\RLN{368.30pt}{والشهور والسنين من قبل ذلك}}}
\At{base east}{522.03pt,-600.64pt}{\hypertarget{AB01-PDF1149-L18}{}{\fontsize{13.8}{27}\ArBody\RLJ{399.30pt}{في معرفة اختلاف حركة الشمس وما يظهر معه من مواضع بعدها الأبعد من اجزاء}}}
\At{base east}{548.43pt,-600.64pt}{{\fontsize{13.8}{27}\ArBody\Ov{كح}}}
\At{base east}{520.53pt,-627.64pt}{\hypertarget{AB01-PDF1149-L19}{}{\fontsize{13.8}{27}\ArBody\RLN{23.70pt}{البروج}}}
\At{base east}{522.03pt,-654.35pt}{\hypertarget{AB01-PDF1149-L20}{}{\fontsize{13.8}{27}\ArBody\RLJ{399.60pt}{في معرفة اقدار اختلاف الايّام بلياليها اذا قيس نهار يوم مع ليلته الى نهار يوم آخر مع ليلته وكيف}}}
\At{base east}{547.73pt,-654.35pt}{{\fontsize{13.8}{27}\ArBody\Ov{كط}}}
\At{base west}{555.33pt,-654.35pt}{{\fontsize{8}{9}\selectfont 20}}
\At{base east}{522.03pt,-681.35pt}{\hypertarget{AB01-PDF1149-L21}{}{\fontsize{13.8}{27}\ArBody\RLN{141.10pt}{تحوَّل وتُنْقَل من بعضها الى بعض}}}
\At{base east}{522.03pt,-708.35pt}{\hypertarget{AB01-PDF1149-L22}{}{\fontsize{13.8}{27}\ArBody\RLJ{399.60pt}{في صفة افلاك القمر وحركاته وما يظهر فيها من الاختلاف في اوقات الاجتماعات والمقابلات}}}
\At{base east}{548.43pt,-708.35pt}{{\fontsize{13.8}{27}\ArBody\Ov{ل}}}
\At{west}{122.23pt,-737.45pt}{\rule{102.2pt}{0.5pt}}
\At{base west}{150.73pt,-764.85pt}{\hypertarget{AB01-PDF1149-N01}{}{\fontsize{9.2}{11}\selectfont 1) Deest in codice. — 2) Cod. \ArRun{زمان}}}
\DSource{1148}{PAG. 4}
\At{base east}{489.93pt,-137.51pt}{\hypertarget{AB01-PDF1148-L01}{}{\fontsize{13.8}{27}\ArBody\RLJ{397.40pt}{الشمسيّة وما يتركّب مع ذلك من الاختلاف الثاني من قبل ابعاده عن الشمس *وعلل الكسوفين}}}
\At{base west}{521.83pt,-137.51pt}{{\fontsize{8}{9}\selectfont f. 3,r.}}
\At{base east}{491.43pt,-165.51pt}{\hypertarget{AB01-PDF1148-L02}{}{\fontsize{13.8}{27}\ArBody\RLN{296.70pt}{وبعد النيّريْن عن الارض وزيادة ضوء القمر ونقصانه ببعده عن الشمس}}}
\At{base east}{491.03pt,-191.11pt}{\hypertarget{AB01-PDF1148-L03}{}{\fontsize{13.8}{27}\ArBody\RLN{173.50pt}{في صفة افلاك الكواكب المتحيّرة وحالاتها}}}
\At{base east}{515.73pt,-191.11pt}{{\fontsize{13.8}{27}\ArBody\Ov{لا}}}
\At{base east}{491.03pt,-218.81pt}{\hypertarget{AB01-PDF1148-L04}{}{\fontsize{13.8}{27}\ArBody\RLN{325.40pt}{في معرفة تأريخ العرب والروم والفرس والقبط ومعرفة بعض ذلك من بعض}}}
\At{base east}{515.73pt,-218.81pt}{{\fontsize{13.8}{27}\ArBody\Ov{لب}}}
\At{base east}{492.13pt,-244.71pt}{\hypertarget{AB01-PDF1148-L05}{}{\fontsize{13.8}{27}\ArBody\RLN{376.60pt}{في معرفة موضع الشمس الذي تُرَى فيه من فلك البروج بتأريخ الروم والعرب ايّهما شئتَ}}}
\At{base east}{516.43pt,-244.71pt}{{\fontsize{13.8}{27}\ArBody\Ov{لج}}}
\At{base east}{82.63pt,-244.71pt}{{\fontsize{8}{9}\selectfont 5}}
\At{base east}{492.53pt,-272.81pt}{\hypertarget{AB01-PDF1148-L06}{}{\fontsize{13.8}{27}\ArBody\RLJ{400.00pt}{في معرفة ساعات التقويم في كلّ بلد وهي الساعات المعتدلة الوسطى التي تكون من بعد انتصاف}}}
\At{base east}{516.03pt,-272.81pt}{{\fontsize{13.8}{27}\ArBody\Ov{لد}}}
\At{base east}{492.53pt,-299.11pt}{\hypertarget{AB01-PDF1148-L07}{}{\fontsize{13.8}{27}\ArBody\RLJ{400.00pt}{النهار بمدينة الرَّقَّة وبها تُستخرَج الحركات في كلّ حين فيُعْرَف وسط الكوكب في ذلك الوقت}}}
\At{base east}{492.53pt,-326.51pt}{\hypertarget{AB01-PDF1148-L08}{}{\fontsize{13.8}{27}\ArBody\RLN{289.10pt}{من اوقات النهار والليل وتحويل هذه الساعات الى ساعات البلدان}}}
\At{base east}{492.53pt,-353.11pt}{\hypertarget{AB01-PDF1148-L09}{}{\fontsize{13.8}{27}\ArBody\RLJ{400.00pt}{في إقامة الطالع والبيوت الاثناعشر من قبل ساعات النهار والليل ومعرفة الساعات من قبل الطالع}}}
\At{base east}{517.13pt,-353.11pt}{{\fontsize{13.8}{27}\ArBody\Ov{له}}}
\At{base east}{492.53pt,-379.41pt}{\hypertarget{AB01-PDF1148-L10}{}{\fontsize{13.8}{27}\ArBody\RLJ{400.00pt}{في معرفة موضع القمر الحقيقيّ من فلك البروج في كلّ يوم وفي كلّ وقت}}}
\At{base east}{517.13pt,-379.41pt}{{\fontsize{13.8}{27}\ArBody\Ov{لو}}}
\At{base east}{84.03pt,-379.41pt}{{\fontsize{8}{9}\selectfont 10}}
\At{base east}{492.53pt,-406.41pt}{\hypertarget{AB01-PDF1148-L11}{}{\fontsize{13.8}{27}\ArBody\RLJ{398.60pt}{في معرفة موضع العقد الشماليّ والجنوبيّ وهما الرأس والذنب اللذين يكون عليهما مجاز القمر}}}
\At{base east}{517.83pt,-406.41pt}{{\fontsize{13.8}{27}\ArBody\Ov{لز}}}
\At{base east}{492.53pt,-433.71pt}{\hypertarget{AB01-PDF1148-L12}{}{\fontsize{13.8}{27}\ArBody\RLN{275.10pt}{في العرض}}}
\At{base east}{492.53pt,-461.11pt}{\hypertarget{AB01-PDF1148-L13}{}{\fontsize{13.8}{27}\ArBody\RLN{321.90pt}{في معرفة عرض القمر وهو بعده عن نطاق البروج الى جهة الجنوب والشمال}}}
\At{base east}{517.83pt,-461.11pt}{{\fontsize{13.8}{27}\ArBody\Ov{لح}}}
\At{base east}{492.53pt,-488.11pt}{\hypertarget{AB01-PDF1148-L14}{}{\fontsize{13.8}{27}\ArBody\RLJ{400.00pt}{في معرفة اختلاف المنظر الذي يعرض في القمر في الطول والعرض واقداره في نواحي الافق}}}
\At{base east}{517.83pt,-488.11pt}{{\fontsize{13.8}{27}\ArBody\Ov{لط}}}
\At{base east}{492.53pt,-514.01pt}{\hypertarget{AB01-PDF1148-L15}{}{\fontsize{13.8}{27}\ArBody\RLN{202.70pt}{والسبب الذي يعرض عنه ذلك فيه بجهات شتَّى}}}
\At{base east}{85.13pt,-514.01pt}{{\fontsize{8}{9}\selectfont 15}}
\At{base east}{492.53pt,-541.41pt}{\hypertarget{AB01-PDF1148-L16}{}{\fontsize{13.8}{27}\ArBody\RLJ{400.00pt}{في معرفة بعد القمر عن الارض من قبل اختلاف منظره في دائرة الارتفاع التي فيما بين سمت}}}
\At{base east}{518.93pt,-541.41pt}{{\fontsize{13.8}{27}\ArBody\Ov{م}}}
\At{base east}{491.73pt,-568.71pt}{\hypertarget{AB01-PDF1148-L17}{}{\fontsize{13.8}{27}\ArBody\RLN{216.30pt}{الرؤس والافق القاطعة لموضع القمر من فلك البروج}}}
\At{base east}{492.53pt,-597.21pt}{\hypertarget{AB01-PDF1148-L18}{}{\fontsize{13.8}{27}\ArBody\RLJ{398.90pt}{في معرفة رؤية الهلال\Note{1} في اوائل الشهور واواخرها وسمت موضعه الذي يُرَى فيه من السماء}}}
\At{base east}{518.93pt,-597.21pt}{{\fontsize{13.8}{27}\ArBody\Ov{ما}}}
\At{base east}{492.53pt,-621.71pt}{\hypertarget{AB01-PDF1148-L19}{}{\fontsize{13.8}{27}\ArBody\RLN{360.00pt}{وارتفاعه عند ذلك عن الافق وصورته على ما *فيه من الضوء واعتدال طرفيْه او ميلهما}}}
\At{base west}{526.13pt,-621.71pt}{{\fontsize{8}{9}\selectfont f. 3,v.}}
\At{base east}{492.53pt,-648.71pt}{\hypertarget{AB01-PDF1148-L20}{}{\fontsize{13.8}{27}\ArBody\RLJ{398.90pt}{في معرفة اجتماعات ومقابلات الشمس والقمر الوسطى والحقيقيّة بتأريخ الروم والقبط ايّهما شئت}}}
\At{base east}{518.93pt,-648.71pt}{{\fontsize{13.8}{27}\ArBody\Ov{مب}}}
\At{base east}{86.23pt,-648.71pt}{{\fontsize{8}{9}\selectfont 20}}
\At{base east}{492.53pt,-675.71pt}{\hypertarget{AB01-PDF1148-L21}{}{\fontsize{13.8}{27}\ArBody\RLJ{397.80pt}{في معرفة الكسوفات القمريّة واقدارها واوقاتها في البلدان والناحية التي منها يبتدئ الكسوف}}}
\At{base east}{519.33pt,-675.71pt}{{\fontsize{13.8}{27}\ArBody\Ov{مج}}}
\At{base east}{492.53pt,-704.11pt}{\hypertarget{AB01-PDF1148-L22}{}{\fontsize{13.8}{27}\ArBody\RLN{363.30pt}{والناحية التي منها يكون الانجلاء من دائرة القمر وصورة ذلك وعمله بالحساب والجدول}}}
\At{west}{93.43pt,-734.71pt}{\rule{101.9pt}{0.5pt}}
\At{base west}{123.73pt,-762.81pt}{\hypertarget{AB01-PDF1148-N01}{}{\fontsize{9.2}{11}\selectfont 1) Cod. \ArRun{الاهلة}}}
\DSource{1147}{PAG. 5}
\At{base east}{515.13pt,-143.83pt}{\hypertarget{AB01-PDF1147-L01}{}{\fontsize{13.8}{27}\ArBody\RLJ{398.80pt}{في معرفة كسوف الشمس واقداره المختلفة في كلّ بلد واوقاته فيه ومعرفة الناحية التي منها يبتدئ}}}
\At{base east}{536.23pt,-143.83pt}{{\fontsize{13.8}{27}\ArBody\Ov{مد}}}
\At{base east}{515.93pt,-170.83pt}{\hypertarget{AB01-PDF1147-L02}{}{\fontsize{13.8}{27}\ArBody\RLN{307.50pt}{وينجلى الكسوف من دائرة الشمس وصورة ذلك وعمله بالحساب والجدول}}}
\At{base east}{515.93pt,-197.43pt}{\hypertarget{AB01-PDF1147-L03}{}{\fontsize{13.8}{27}\ArBody\RLN{306.80pt}{في معرفة مواضع الخمسة الكواكب المتحيّرة من فلك البروج في كلّ حين}}}
\At{base east}{538.03pt,-197.43pt}{{\fontsize{13.8}{27}\ArBody\Ov{مه}}}
\At{base east}{515.93pt,-224.43pt}{\hypertarget{AB01-PDF1147-L04}{}{\fontsize{13.8}{27}\ArBody\RLN{206.30pt}{في معرفة مقام الكواكب الخمسة المتحيّرة ورجوعها}}}
\At{base east}{538.33pt,-224.43pt}{{\fontsize{13.8}{27}\ArBody\Ov{مو}}}
\At{base east}{515.93pt,-251.13pt}{\hypertarget{AB01-PDF1147-L05}{}{\fontsize{13.8}{27}\ArBody\RLN{215.70pt}{في معرفة عروض الكواكب الخمسة المتحيّرة وجهاتها}}}
\At{base east}{537.33pt,-251.13pt}{{\fontsize{13.8}{27}\ArBody\Ov{مز}}}
\At{base west}{544.53pt,-251.13pt}{{\fontsize{8}{9}\selectfont 5}}
\At{base east}{515.93pt,-278.13pt}{\hypertarget{AB01-PDF1147-L06}{}{\fontsize{13.8}{27}\ArBody\RLN{249.50pt}{في معرفة ظهور الكواكب الخمسة المتحيّرة واختفائها}}}
\At{base east}{539.13pt,-278.13pt}{{\fontsize{13.8}{27}\ArBody\Ov{مح}}}
\At{base east}{515.93pt,-305.83pt}{\hypertarget{AB01-PDF1147-L07}{}{\fontsize{13.8}{27}\ArBody\RLN{348.50pt}{في معرفة الاشكال التسعة التي تكون للكواكب الثابتة وبعض المتحيّرة عند الشمس}}}
\At{base east}{539.43pt,-305.83pt}{{\fontsize{13.8}{27}\ArBody\Ov{مط}}}
\At{base east}{515.93pt,-332.83pt}{\hypertarget{AB01-PDF1147-L08}{}{\fontsize{13.8}{27}\ArBody\RLN{313.20pt}{في ذكر ابعاد الكواكب عن الارض واقطارها وعظم اجرامها وسعة افلاكها}}}
\At{base east}{539.83pt,-332.83pt}{{\fontsize{13.8}{27}\ArBody\Ov{ن}}}
\At{base east}{515.93pt,-359.43pt}{\hypertarget{AB01-PDF1147-L09}{}{\fontsize{13.8}{27}\ArBody\RLJ{399.60pt}{في معرفة حركة سائر الكواكب\Note{1} بالرصد ورسم مواضع ما يحتاج اليه منها في الجدول في الطول}}}
\At{base east}{540.13pt,-359.43pt}{{\fontsize{13.8}{27}\ArBody\Ov{نا}}}
\At{base east}{515.93pt,-386.43pt}{\hypertarget{AB01-PDF1147-L10}{}{\fontsize{13.8}{27}\ArBody\RLN{41.80pt}{والعرض}}}
\At{base west}{547.33pt,-386.43pt}{{\fontsize{8}{9}\selectfont 10}}
\At{base east}{515.93pt,-413.43pt}{\hypertarget{AB01-PDF1147-L11}{}{\fontsize{13.8}{27}\ArBody\RLN{374.80pt}{فيما ذكر اصحاب الطلسمات انّ للفلك حركة انتقال مقبلةً ومدبرةً وما يظهر فيه من الخلل}}}
\At{base east}{540.13pt,-413.43pt}{{\fontsize{13.8}{27}\ArBody\Ov{نب}}}
\At{base east}{515.93pt,-440.83pt}{\hypertarget{AB01-PDF1147-L12}{}{\fontsize{13.8}{27}\ArBody\RLJ{399.60pt}{في معرفة اوقات تحاويل السنين الكائنة\Note{2} عند عودة الشمس الى الموضع الذي كانت فيه في الاصل}}}
\At{base east}{541.23pt,-440.83pt}{{\fontsize{13.8}{27}\ArBody\Ov{نج}}}
\At{base east}{515.93pt,-467.83pt}{\hypertarget{AB01-PDF1147-L13}{}{\fontsize{13.8}{27}\ArBody\RLJ{399.60pt}{في تحقيق اقدار الاتّصالات التي تكون بحسب عروض الكواكب اذا ألقت الشعاع على فلك}}}
\At{base east}{541.63pt,-467.83pt}{{\fontsize{13.8}{27}\ArBody\Ov{ند}}}
\At{base east}{515.13pt,-494.43pt}{\hypertarget{AB01-PDF1147-L14}{}{\fontsize{13.8}{27}\ArBody\RLN{24.40pt}{البروج}}}
\At{base east}{515.93pt,-520.73pt}{\hypertarget{AB01-PDF1147-L15}{}{\fontsize{13.8}{27}\ArBody\RLN{185.80pt}{في معرفة مطالع البروج فيما بين ارباع الفلك}}}
\At{base east}{541.63pt,-520.73pt}{{\fontsize{13.8}{27}\ArBody\Ov{نه}}}
\At{base west}{548.83pt,-520.73pt}{{\fontsize{8}{9}\selectfont 15}}
\At{base east}{515.93pt,-548.12pt}{\hypertarget{AB01-PDF1147-L16}{}{\fontsize{13.8}{27}\ArBody\RLJ{399.60pt}{في عمل *الرخامة القائمة المسطوحة لمعرفة ساعات النهار الزمانيّة في كلّ بلد وتقويم نصبها وسمت}}}
\At{base east}{541.93pt,-548.12pt}{{\fontsize{13.8}{27}\ArBody\Ov{نو}}}
\At{base east}{108.13pt,-548.12pt}{{\fontsize{8}{9}\selectfont f. 4,r.}}
\At{base east}{515.13pt,-575.12pt}{\hypertarget{AB01-PDF1147-L17}{}{\fontsize{13.8}{27}\ArBody\RLN{300.90pt}{الجنوب وكيف يُعْرَف سمت القبلة في الرخامة وهو سمت مكّة المحروسة}}}
\At{base east}{515.93pt,-602.12pt}{\hypertarget{AB01-PDF1147-L18}{}{\fontsize{13.8}{27}\ArBody\RLJ{399.60pt}{في ختم الكتاب وصفة ضعة الآلة التي على هيئة الفلك وتسمَّى البيْضة وضعة الآلتين اللّتان\Note{3}}}}
\At{base east}{542.73pt,-602.12pt}{{\fontsize{13.8}{27}\ArBody\Ov{نز}}}
\At{base east}{515.93pt,-628.73pt}{\hypertarget{AB01-PDF1147-L19}{}{\fontsize{13.8}{27}\ArBody\RLN{37.10pt}{للرصد.}}}
\At{base east}{515.93pt,-655.42pt}{\hypertarget{AB01-PDF1147-L20}{}{\fontsize{13.8}{27}\ArBody\RLN{380.20pt}{وهذا تفسير تفصيل الكتاب وهو سبعة وخمسون نوعًا والحمد لله على عونه وصلّى الله على محمّد.}}}
\At{base west}{521.13pt,-655.42pt}{{\fontsize{8}{9}\selectfont 20}}
\At{west}{116.13pt,-724.12pt}{\rule{101.5pt}{0.5pt}}
\At{base west}{144.63pt,-751.53pt}{\hypertarget{AB01-PDF1147-N01}{}{\fontsize{9.2}{11}\selectfont \makebox[373.30pt][s]{1 Pro \ArRun{سائر الكواكب}, in capite ipso et apud Platonem legitur \ArRun{الكواكب الثابتة} — 2 Cod. \ArRun{الكائن} —}}}
\At{base west}{118.33pt,-767.03pt}{\hypertarget{AB01-PDF1147-N02}{}{\fontsize{9.2}{11}\selectfont 3) Cod. \ArRun{التي}}}
\DSource{1146}{PAG. 6}
\At{base}{302.73pt,-160.02pt}{\hypertarget{AB01-PDF1146-L01}{}{\fontsize{21.0}{25.2}\ArBody\RL{الباب الاوّل}}}
\At{base}{301.78pt,-203.62pt}{\hypertarget{AB01-PDF1146-L02}{}{\fontsize{15.5}{18.6}\ArBody\RL{في صدر الكتاب}}}
\At{base east}{491.73pt,-250.02pt}{\hypertarget{AB01-PDF1146-L03}{}{\fontsize{13.8}{27}\ArBody\RLJ{401.70pt}{قال إِنّ اوّلَ ما أُبْتُدِئَ به كلُّ امر وأُسْتُفْتِحَ به كلُّ قول حمد الله جلّ ذكره والثناء عليه بآلائه}}}
\At{base east}{81.93pt,-250.02pt}{{\fontsize{8}{9}\selectfont 5}}
\At{base east}{512.33pt,-277.43pt}{\hypertarget{AB01-PDF1146-L04}{}{\fontsize{13.8}{27}\ArBody\RLJ{422.30pt}{والصلاة على خاتم رُسُله وأنبيائه عليهم السلام ورحمة الله وبركاته. ﴿ الحمد لله الذي خلق الخلائق}}}
\At{base east}{512.33pt,-304.02pt}{\hypertarget{AB01-PDF1146-L05}{}{\fontsize{13.8}{27}\ArBody\RLJ{422.30pt}{بقدرته ودبّر الامور بمشيئته وأتقنها بحكمته ﴾ فأَحاطَ بِكُلِّ شَيْءٍ عِلْمًا\Note{1} وَأَحْصَى كُلَّ شَيْءٍ عَدَدًا\Note{2} لا}}}
\At{base east}{512.33pt,-330.62pt}{\hypertarget{AB01-PDF1146-L06}{}{\fontsize{13.8}{27}\ArBody\RLJ{422.30pt}{يَعْزُبُ عَنْهُ مِثْقالُ ذَرَّةٍ في ٱلسَّماواتِ وَلا في ٱلْأَرْضِ وَلا أَصْغَرُ مِنْ ذَلِكَ وَلا أَكْبَرُ إِلَّا في كِتَابٍ}}}
\At{base east}{512.33pt,-357.62pt}{\hypertarget{AB01-PDF1146-L07}{}{\fontsize{13.8}{27}\ArBody\RLJ{422.30pt}{مُبِينٍ\Note{3} واشهدُ أَن لا اله إِلَّا الله وحْدَه لا شريكَ له واشهدُ أَنَّ محمّدًا عبده ورسوله أَرْسَلَهُ بِٱلهُدَى}}}
\At{base east}{512.33pt,-385.02pt}{\hypertarget{AB01-PDF1146-L08}{}{\fontsize{13.8}{27}\ArBody\RLJ{422.30pt}{وَدِينِ ٱلْحَقِّ لِيُظْهِرَهُ عَلَى ٱلدِّينِ كُلِّهِ ولَوْ كَرِهَ ٱلْمُشْرِكُونَ\Note{4} فهدى به المؤمنين وقَطَعَ بِهِ دَابِرَ ٱلكَافِرِينَ\Note{5}}}}
\At{base east}{82.23pt,-385.02pt}{{\fontsize{8}{9}\selectfont 10}}
\At{base east}{512.33pt,-412.02pt}{\hypertarget{AB01-PDF1146-L09}{}{\fontsize{13.8}{27}\ArBody\RLJ{422.30pt}{وجعله حجّة على العالمين صلّى الله عليه وعلى آله الطيّبين وعلى اصحابه المنتخبين وعلى التابعين لسنّته الى}}}
\At{base east}{512.33pt,-439.02pt}{\hypertarget{AB01-PDF1146-L10}{}{\fontsize{13.8}{27}\ArBody\RLJ{422.30pt}{يوم الدين. ﴿ امّا بعد ﴾ إنّ من اشرف العلوم منزلةً واسناها مرتبةً واحسنها حِليةً واعلقها بالقلوب}}}
\At{base east}{512.33pt,-466.43pt}{\hypertarget{AB01-PDF1146-L11}{}{\fontsize{13.8}{27}\ArBody\RLJ{422.30pt}{والمعها بالنفوس واشدّها تحديدًا للفكر والنظر وتذكيةً للفهْم ورياضةً للعقل بعد العلم بما لا يسع الانسانَ}}}
\At{base east}{512.33pt,-493.02pt}{\hypertarget{AB01-PDF1146-L12}{}{\fontsize{13.8}{27}\ArBody\RLJ{422.30pt}{جهلُه من شرائع الدين وسنّته علم صناعة النجوم لِمَا في ذلك من جسيم الحظّ وعظيم الانتفاع بمعرفة}}}
\At{base east}{512.33pt,-520.42pt}{\hypertarget{AB01-PDF1146-L13}{}{\fontsize{13.8}{27}\ArBody\RLJ{422.30pt}{مُدّة السنين والشهور *والمواقيت وفصول الازمان وزيادة النهار واللّيل ونقصانها ومواضع النيّريْن}}}
\At{base east}{82.63pt,-520.42pt}{{\fontsize{8}{9}\selectfont 15}}
\At{base west}{521.43pt,-520.42pt}{{\fontsize{8}{9}\selectfont f. 4,v.}}
\At{base east}{512.33pt,-547.42pt}{\hypertarget{AB01-PDF1146-L14}{}{\fontsize{13.8}{27}\ArBody\RLJ{422.30pt}{وكسوفهما ومسير الكواكب في استقامتها ورجوعها وتبدّل اشكالها ومراتب افلاكها وسائر مناسَباتها الى}}}
\At{base east}{512.33pt,-574.03pt}{\hypertarget{AB01-PDF1146-L15}{}{\fontsize{13.8}{27}\ArBody\RLJ{422.30pt}{ما يُدْرِك بذلك مَنْ أنعم النظر وأدام الفكر فيه من إثبات التوحيد ومعرفة كُنْه عظمة الخالق وسَعة}}}
\At{base east}{512.33pt,-601.03pt}{\hypertarget{AB01-PDF1146-L16}{}{\fontsize{13.8}{27}\ArBody\RLJ{422.30pt}{حكمته وجليل قدرته ولطيف صنعه قال عزّمَن قائل\Note{6} إِنَّ في خَلْقِ ٱلسَّماوَاتِ وَٱلْأَرْضِ وَٱخْتِلَافِ ٱللَّيْلِ}}}
\At{base east}{512.33pt,-628.03pt}{\hypertarget{AB01-PDF1146-L17}{}{\fontsize{13.8}{27}\ArBody\RLJ{422.30pt}{وَٱلنَّهَارِ لَآيَاتٍ لِأُولِي ٱلْأَلْبَابِ\Note{7} وقال تبارك وتعالى تَبَارَكَ ٱلَّذِي جَعَلَ في ٱلْسَّمَآء بُرُوجًا.\Note{8} وقال عزّ وجلّ}}}
\At{base east}{512.33pt,-654.62pt}{\hypertarget{AB01-PDF1146-L18}{}{\fontsize{13.8}{27}\ArBody\RLJ{422.30pt}{هُوَ ٱلَّذِي جَعَلَ ٱللَّيْلَ وَٱلْقَمَرَ نُورًا خِلْفَةً.\Note{9} وقال سبحانه هُوَ ٱلَّذِي جَعَلَ ٱلشَّمْسَ ضِيَاءً وَٱلْقَمَرَ نُورًا}}}
\At{base east}{82.63pt,-654.62pt}{{\fontsize{8}{9}\selectfont 20}}
\At{base east}{512.33pt,-681.33pt}{\hypertarget{AB01-PDF1146-L19}{}{\fontsize{13.8}{27}\ArBody\RLJ{422.30pt}{وَقَدَّرَهُ مَنَازِلَ لِتَعْلَمُوا عَدَدَ ٱلسِّنِينَ وَٱلْحِسَابَ.\Note{10} وقال جلّ ذكره ٱلشَّمْسُ وَٱلْقَمَرُ بِحُسْبَانٍ\Note{11} مع اقتصاص}}}
\At{west}{89.43pt,-709.22pt}{\rule{102.3pt}{0.5pt}}
\At{base west}{117.93pt,-734.62pt}{\hypertarget{AB01-PDF1146-N01}{}{\fontsize{9.2}{11}\selectfont \makebox[396.40pt][s]{1) Qor. LXV, 12. — 2) Qor. LXXII, 28. — 3) Qor. XXXIV, 3. — 4) Qor. IX, 33 et XLI, 9. — 5) Cfr.}}}
\At{base west}{92.03pt,-749.72pt}{\hypertarget{AB01-PDF1146-N02}{}{\fontsize{9.2}{11}\selectfont \makebox[422.30pt][s]{Qor. VII, 70 et VIII, 7. — 6) Cod. \ArRun{قابل}. — 7) Qor. III, 187. — 8) Qor. XXV, 62. — 9) Qor. XXV, 63. —}}}
\At{base west}{92.03pt,-765.22pt}{\hypertarget{AB01-PDF1146-N03}{}{\fontsize{9.2}{11}\selectfont 10) Qor. X, 5. — 11) Qor. LV, 4.}}
\DSource{1145}{PAG. 7}
\At{base east}{548.33pt,-143.84pt}{\hypertarget{AB01-PDF1145-L01}{}{\fontsize{13.8}{27}\ArBody\RLJ{423.00pt}{كثير في كتاب الله عزّ وجلّ يطول وصفه ويتّسع القول بذكره واستشهاده. ﴿ وإنّي\Note{1} لمّا اطَالتُ}}}
\At{base east}{549.03pt,-170.44pt}{\hypertarget{AB01-PDF1145-L02}{}{\fontsize{13.8}{27}\ArBody\RLJ{423.70pt}{النظر ﴾ في هذا العلم وادمنتُ الفكر فيه ووقفتُ على اختلاف الكتب الموضوعة لحركات النجوم وما}}}
\At{base east}{550.13pt,-197.15pt}{\hypertarget{AB01-PDF1145-L03}{}{\fontsize{13.8}{27}\ArBody\RLJ{424.80pt}{تَهَيَّأَ\Note{2} على بعض واضعيها من الخَلَل فيما أَصَّلوه فيها من الاعمال وما ابتنوها\Note{3} عليه وما اجتمع ايضًا في}}}
\At{base east}{550.13pt,-223.75pt}{\hypertarget{AB01-PDF1145-L04}{}{\fontsize{13.8}{27}\ArBody\RLJ{424.80pt}{حركات النجوم على طول الزمان لمّا قيست أرصادها الى الأرصاد القديمة وما وُجِدَ في ميْل فلك البروج}}}
\At{base east}{550.43pt,-250.75pt}{\hypertarget{AB01-PDF1145-L05}{}{\fontsize{13.8}{27}\ArBody\RLJ{425.10pt}{عن فلك معدِّل النهار من التقارب وما تغيّر بتغيّره من اصناف الحساب واقدار ازمان السنين واوقات}}}
\At{base west}{556.73pt,-250.75pt}{{\fontsize{8}{9}\selectfont 5}}
\At{base east}{550.43pt,-278.14pt}{\hypertarget{AB01-PDF1145-L06}{}{\fontsize{13.8}{27}\ArBody\RLJ{425.10pt}{الفصول واتّصالات النيّريْن التي يُستدلّ عليها بازمان الكسوفات واوقاتها اجريتُ في تصحيح ذلك}}}
\At{base east}{551.13pt,-305.44pt}{\hypertarget{AB01-PDF1145-L07}{}{\fontsize{13.8}{27}\ArBody\RLJ{425.50pt}{وإحكامه على مذهب بطلميوس في الكتاب المعروف بالمجسطيّ بعد إنعام النَظَر وطول الفكر والروْية}}}
\At{base east}{551.13pt,-331.75pt}{\hypertarget{AB01-PDF1145-L08}{}{\fontsize{13.8}{27}\ArBody\RLJ{425.10pt}{مقتفيًا اثره متّبعًا ما رسمه اذ كان قد تَقَصَّى ذلك من وجوهه ودلّ على العلل والاسباب العارضة}}}
\At{base east}{551.53pt,-358.75pt}{\hypertarget{AB01-PDF1145-L09}{}{\fontsize{13.8}{27}\ArBody\RLJ{426.20pt}{*فيه بالبرهان الهندسيّ والعدديّ الذي لا تُدْفَع صحّته ولا يُشَكّ في حقيقته فأمر بالمِحْنة والاعتبار}}}
\At{base east}{118.93pt,-358.75pt}{{\fontsize{8}{9}\selectfont f. 5,r.}}
\At{base east}{551.53pt,-385.75pt}{\hypertarget{AB01-PDF1145-L10}{}{\fontsize{13.8}{27}\ArBody\RLJ{425.50pt}{بعده وذكر انّه قد يجوز أن يُسْتَدْرَكَ عليه في أرصاده على طول الزمان كما استدرك هو على إبَّرْخُس\Note{4}}}}
\At{base west}{559.63pt,-385.75pt}{{\fontsize{8}{9}\selectfont 10}}
\At{base east}{551.53pt,-412.75pt}{\hypertarget{AB01-PDF1145-L11}{}{\fontsize{13.8}{27}\ArBody\RLJ{424.80pt}{وغيره من نظرائه لجلالة الصناعة ولأنّها سمائيّة جسيمة لا تُدْرَك الّا بالتقريب ووضعتُ في ذلك}}}
\At{base east}{551.53pt,-439.35pt}{\hypertarget{AB01-PDF1145-L12}{}{\fontsize{13.8}{27}\ArBody\RLJ{425.20pt}{كتابًا اوضحتُ فيه ما ٱستعجم وفتحتُ ما ٱستغلق وبيّنتُ ما أشكل من اصول هذا العلم وشذّ من فروعه}}}
\At{base east}{551.53pt,-466.35pt}{\hypertarget{AB01-PDF1145-L13}{}{\fontsize{13.8}{27}\ArBody\RLJ{424.80pt}{وسهّلتُ به سبيل الهداية لمن يأثَر به ويعمل عليه في صناعة النجوم\Note{5} وصحّحتُ فيه حركات الكواكب}}}
\At{base east}{551.53pt,-493.05pt}{\hypertarget{AB01-PDF1145-L14}{}{\fontsize{13.8}{27}\ArBody\RLJ{424.40pt}{ومواضعها من مِنطَقة فلك البروج على نحو ما وجدتها بالرصد وحساب الكسوفيْن وسائر ما يُحتاج اليه}}}
\At{base east}{551.53pt,-520.04pt}{\hypertarget{AB01-PDF1145-L15}{}{\fontsize{13.8}{27}\ArBody\RLJ{424.10pt}{من الاعمال وأَضفتُ الى ذلك غيره ممّا يُحتاج اليه وجعلتُ ٱستخراج حركات الكواكب فيه من الجداول}}}
\At{base west}{561.03pt,-520.04pt}{{\fontsize{8}{9}\selectfont 15}}
\At{base east}{551.53pt,-547.04pt}{\hypertarget{AB01-PDF1145-L16}{}{\fontsize{13.8}{27}\ArBody\RLJ{424.40pt}{لوقت انتصاف النهار من اليوم الذي يُحْسَب فيه بمدينة الرَّقَّة وبها كان الرصد والامتحان على تحذيق}}}
\At{base east}{551.53pt,-574.04pt}{\hypertarget{AB01-PDF1145-L17}{}{\fontsize{13.8}{27}\ArBody\RLN{413.60pt}{ذلك كلّه إن شاء الله تعالى وبالله التوفيق.}}}
\At{west}{129.43pt,-731.00pt}{\rule{100.4pt}{0.5pt}}
\At{base west}{158.23pt,-748.64pt}{\hypertarget{AB01-PDF1145-N01}{}{\fontsize{9.2}{11}\selectfont \makebox[396.00pt][s]{1) Cod. \ArRun{وان} — 2) Cod. \ArRun{تهي} — 3) Cod. \ArRun{ابتدوها} — 4) Cod. semper \ArRun{برخس}: sed ceteri Arabes ut recepi.}}}
\At{base west}{132.33pt,-763.75pt}{\hypertarget{AB01-PDF1145-N02}{}{\fontsize{9.2}{11}\selectfont Plato: «Abrachis». — 5) Deest in cod.}}
\end{document}
